\documentclass[10pt,a4paper]{article}
\usepackage[margin=24mm]{geometry}
\usepackage{amsmath,amssymb,booktabs,longtable}
\usepackage[T1]{fontenc}
\usepackage{hyperref}
\title{Geometry- and Physics-Guided IR Train Proof Contract}
\author{Proof repair artifact for Immersive Railroading 1.11.0}
\date{2026-08-16}

\newcommand{\Lean}[1]{\texttt{#1}}
\newcommand{\status}{\textsc{not ready}}
\newcommand{\bytecode}{\textsf{bytecode fact}}
\newcommand{\model}{\textsf{certified model assumption}}
\newcommand{\runtime}{\textsf{runtime contract}}
\newcommand{\obligation}{\textsf{future Lua implementation obligation}}

\begin{document}
\maketitle

\begin{center}
\fbox{\parbox{0.88\linewidth}{\centering
This artifact replaces the PID proof.  Its honest status is
\textbf{NOT\_READY}.  Lean proves the exact-rational contract lemmas listed
below, but the complete runtime refinement from Lua, Java floating point,
actual track samples, and the full multi-particle IR plant is still an
implementation obligation.}}
\end{center}

\tableofcontents

\section{Scope and claim discipline}

The proof target is robust front-coupler arrival at a directed station or
signal, within a derived computational tolerance.  The representation pipeline
is deliberately not a PID pipeline:

\[
  \begin{split}
    \text{tracer samples}
    \longrightarrow \text{track-centre cubic splines}
    &\longrightarrow \text{directed routing graph}
    \longrightarrow \text{detected consist profile} \\
    \longrightarrow \text{front-coupler controller}
    \longrightarrow \text{IR plant}.
  \end{split}
\]

The spline model is geometric source of truth.  The graph is a routing
abstraction.  The graph never edits or fits a spline.  A station marker is not
an arbitrary graph vertex, and a raw entity position is not silently treated as
the rail centreline or the front coupler.

Every result in this document has one of these classifications:
\begin{itemize}
  \item \bytecode: directly observed in the Java-8 root class or manifest;
  \item derived from bytecode: an algebraic consequence of those observations;
  \item \model: an explicit assumption checked at the runtime boundary;
  \item \runtime: a required data, timing, or fail-closed contract;
  \item \obligation: a precise future production-Lua change;
  \item unsupported or withdrawn: intentionally not claimed.
\end{itemize}

\section{Provenance}

The authoritative input is the extracted tree under
\texttt{immersive\_railroading/.cache/jar}.  The manifest reports
\texttt{Multi-Release: true}; the root class files have major version 52
(Java 8), while the classes under \texttt{META-INF/versions/16} have major
version 60.  The root Java-8 classes are the target; version-16 classes are not
used silently.  The mod metadata reports Minecraft 1.7.10 and Immersive
Railroading 1.11.0.

The exact manifest, class hashes, and decompiler-input hashes are recorded in
\texttt{pid-jar-manifest.md} and
\texttt{pid-bytecode-traceability.md}.  The proof does not use a historical
PID theorem as authority.  The old files were inspected only after this model
and bytecode pass, to identify the migration surface and obsolete claims.

\section{Geometry: traced centreline to spline}

\subsection{Cubic representation}

A physical section is a 3D cubic Bezier curve
\[
 B(t)=(1-t)^3p_1+3(1-t)^2tp_2+3(1-t)t^2p_3+t^3p_4,
 \qquad 0\leq t\leq 1.
\]
The Lean structure \Lean{Cubic} stores the four 3D points and
\Lean{cubicPosition} and \Lean{cubicDerivative} use the Bernstein
expressions.  The reverse orientation is the control-point reversal
\Lean{cubicReverse}.  The endpoint and derivative facts are proved by Lean.
The derivative endpoint is $3(p_2-p_1)$ and $3(p_4-p_3)$, not an unscaled
chord.

Straight track is encoded as a degenerate cubic with both interior controls at
the endpoint displacement.  This is a representation choice, not a claim
that all IR curves are straight.  The Java root implementation of
\texttt{CubicCurve.circle()} uses a cubic approximation constant
\texttt{0.55191502449}; therefore IR turns are not described as exact
circles.  Arc length is stored by certified marks or a certified numerical
approximation.  Each spline stores tangent, grade, curvature, fitting error,
trace identifiers, optional builder metadata, and a stable identifier.

The geometric certificate requires monotone $(t,s)$ arc marks, positive total
length, nonnegative fit error, correct reverse controls, and tangent metadata.
It is implemented by \Lean{SplineCertificate}.  Curvature and grade are
metadata/derived quantities and must be bounded before they enter a stopping
certificate.

\subsection{Tracer reference and residual}

The tracer record contains tick, raw IR position, speed, heading, tracer
definition ID, stock UUID, validity, and sample age.  A tracer configuration
contains a reference offset, sampling spacing, curvature bound, floating-point
error, and missing-observation bound.  The centre estimate is
\[
  c_i = r_i + o_{\mathrm{tracer}},
\]
and the certified residual budget is
\[
  \varepsilon_{\mathrm{trace}}
  = \Delta\ell\,\kappa_{\max}
    +\varepsilon_{\mathrm{float}}
    +\varepsilon_{\mathrm{missing}}.
\]
The proof record additionally requires a residual vector $e_i$ with
\[
  c_i^{\mathrm{actual}}=c_i+e_i,
  \qquad \|e_i\|_1\leq\varepsilon_{\mathrm{trace}}.
\]
This is a certified model assumption until the Lua sampling loop and the exact
stock/bogey offsets are implemented and measured.  It explicitly avoids the
invalid claims that raw \texttt{getPos()} is the F3 hitbox or the exact rail
centreline, and that the tracer reference is the front coupler.

Repeated traces are matched in 3D using distance, tangent, curvature, and arc
coordinate.  A match can extend, continue, overlap, split, intersect, or
create new geometry.  Two projected paths that cross in 2D are not merged.
The automatic topology classifier has the following conservative policy:

\begin{center}
\begin{tabular}{ll}
\toprule
Observation & Routing result \\
\midrule
same interval & continuation \\
interval extension & extension \\
validated physical junction & shared connection \\
both paths continue at same level without connection & isolated crossing \\
3D height separation above tolerance & overpass \\
insufficient or contradictory observation & unresolved \\
\bottomrule
\end{tabular}
\end{center}

Only continuation, extension, and validated shared connections are routable.
Crossings preserve independent spline paths.  An overpass is separate when its
3D separation exceeds the certified tolerance.  Unexplored branches,
coincident same-level paths with indistinguishable connectivity, and hidden
runtime connectivity remain unresolved and unavailable for routing; manual
classification is required before activation.

\section{Graph, markers, and route locality}

The routing graph stores spline IDs and arc-length intervals on edges.  Vertices
come from spline endpoints, validated junctions, approved switch/merger nodes,
and marker mappings.  A physical spline normally produces symmetric forward
and reverse edges.  Direction restrictions are imposed on the relevant edge or
interval, not by changing the spline geometry.

A station or signal is mapped to the closest spline only when its 3D distance
is at most its explicit tolerance.  The mapping stores marker ID, spline ID,
arc coordinate, measured distance, kind, and exactly one configured approach
orientation.  A marker too far from every spline is rejected.  The opposite
approach is not accepted unless separately configured.

An active route stores graph revision, spline IDs, vertex IDs, and edge IDs.
Validation checks only those referenced objects and the current graph revision;
it does not scan the entire graph.  A revision change invalidates the route,
invokes route finding against the latest graph, and enters fail-closed stop with
an in-game report when no route exists.  Geometry reconstruction and dynamic
occupancy/routing are separate operations.

\section{Detector profile}

The detector contract consumes the \texttt{ir\_train\_overhead} event and
records its \texttt{stock\_uuid}.  The detector's \texttt{info()} call is
immediate and records definition ID, weight, direction, speed, brake channels,
brake-system and adhesion data, and locomotive data when present.  Catalogue
order is the event order, including repeated definition IDs.  \texttt{consist()}
is used only as aggregate validation; it does not replace per-stock records.

The initial catalogue speed is no greater than approximately 10 km/h.  A
higher speed is permitted only after a detector-length, minimum-spacing, event
latency, and \texttt{info()} response-time bound is measured.  The catalogue
profile is immutable during a run.  A UUID/order/definition/weight or brake
mismatch reports an in-game error, stops the controller, and requires player
recataloguing.

\section{Directed stopping coordinate}

For the selected directed spline orientation, $s$ increases along the active
route.  The front coupler, derived from stock definition and IR coupler
semantics, is the stopping reference:
\[
 x=s_{\mathrm{target}}-s_{\mathrm{front\ coupler}},
 \qquad
 v_s=\frac{d s_{\mathrm{front\ coupler}}}{dt}.
\]
The theorem is only entered under $x>0$ and $v_s>0$.  Wrong-way motion,
away-motion, target crossing, unavailable orientation, invalid localization,
or stale data enters recovery or fail-closed state.  World direction is handled
by selecting the appropriate directed spline orientation; it is not handled by
silently accepting an incorrectly oriented marker.

\section{IR physics derivation}

This section records the equation used by Lean and the controller contract.
The evidence and exact bytecode locations are in the traceability file.

\subsection{Material, efficiency, and pressure branches}

The Java-8 root constants are:
\[
 \mu_{ss}^{static}=0.70=\frac{7}{10},\quad
 \mu_{ss}^{kinetic}=0.42=\frac{21}{50},\quad
 \mu_{sci}^{kinetic}=0.25=\frac14,
 \quad g=9.8=\frac{49}{5}.
\]
The effective pressure is the maximum of train-brake and independent-brake
position, normalized to the configured range.  In the contract model it is
\[
 p=\min(1,\max(p_{train},p_{ind})).
\]
This explicitly retains both brake channels.  Normalization rejects NaN and
clamps values outside $[-1,1]$ at the API boundary; invalid or unavailable
values are a runtime failure, not a mathematical zero.

For a physics input record $z$, the ordinary brake demand is
\[
 F_{brake,design}
   =m_{design}g\mu_{sci}^{kinetic}\,p,
\]
then the wheel-slip branch replaces it with the steel kinetic term
\[
 F_{brake,slip}=m g\mu_{ss}^{kinetic}
\]
when the root locomotive branch reports slipping and $|v|>0.01$ m/s.  The
configured \texttt{brakeMultiplier} multiplies the selected branch.  The
maximum adhesion term uses current mass, steel static friction, gravity, and
the separate brake-adhesion efficiency.  Cogging locomotives override both
brake-system and brake-adhesion efficiencies to 10 in the inspected root
override; ordinary rolling stock uses its configured brake-shoe friction and
base adhesion efficiency.

\subsection{Signed equation}

The signed route force is derived in this order, preserving signs and units:
\begin{align*}
 F_{grade} &= m(-g)\sin(\theta)\,M_{slope},\\
 F_{traction} &= u_{throttle}F_{available},\\
 F_{rdi} &= F_{rolling}+F_{direct}+F_{interference}+F_{near-zero},\\
 F_{coupler} &= F_{curvature}+F_{slack}+F_{push/pull},\\
 F_{net} &= F_{grade}+F_{traction}-F_{rdi}-F_{brake}-F_{coupler},\\
 a &= F_{net}/m.
\end{align*}
The Lean definitions \Lean{gradeForceN}, \Lean{tractiveForceN},
\Lean{rollingDirectInterferenceN}, \Lean{couplerAdverseN},
\Lean{brakeForceN}, and \Lean{signedNetForceN} are connected by the
proved \Lean{signed\_force\_equation} identity.  The force ledger is not an
unexplained abstraction: each summand corresponds to one root bytecode branch.

The root \texttt{SimulationState.forcesNewtons()} contributes grade and
tractive force.  \texttt{frictionNewtons()} contributes rolling resistance,
near-zero resistance, interference resistance, brake candidate, the adhesion
wheel-slip substitution, the configured multiplier, and direct resistance.
The root \texttt{Consist} particle path applies signed force and friction per
particle, then aggregates mass, pressure, linkages, slack, and push/pull state
over discrete steps.  The Lean \Lean{ParticleState}, \Lean{ConsistState},
\Lean{irParticleStep}, and \Lean{consistStep} expose this boundary; the
full Java linkage/collision implementation remains a refinement obligation.

\section{Combined command and plant transition}

The controller command is
\[
 u=(u_{throttle},u_{train\ brake},u_{independent\ brake},u_{emergency}).
\]
The stopping command deliberately allows a positive throttle cap in the model;
it is not a hidden zero-throttle assumption.  Its upper positive traction is
included in the stopping bound.  The transition is:
\[
 \begin{split}
  \text{spline localization}
  \to x,v_s
  \to (m,L,\kappa,\text{slack},\text{push/pull})
  \to u_{throttle}
  \to (u_{train\ brake},u_{independent\ brake}) \\
  \to F_{net}
  \to \text{particle step}
  \to \text{next controller state}.
 \end{split}
\]
\Lean{controllerToCombinedForce} (rendered as
\Lean{controller\_to\_combined\_force}) expands the five signed force
branches under the command update.  \Lean{plantTransition} then writes the
same command-modified physics record into the next particle state.  The
\Lean{plant\_transition\_connects\_command\_and\_force} theorem and
\Lean{plant\_braking\_upper} connect command, force, acceleration, and
discrete velocity.  Saturation is represented by \Lean{clamp}; slew limits,
pressure propagation, and emergency-channel timing are explicit runtime
parameters still required from Lua.

\section{Uniform guard and safety envelope}

Let $B_{min}$ be a lower brake-force bound valid for every particle and every
branch over the complete observation-plus-actuation-plus-pressure delay
horizon.  Let $T_{max}$ bound positive traction and $A_{max}$ bound adverse
grade, curvature, slack, push/pull, direct-resistance uncertainty, and
localization uncertainty in force units.  The uniform positive deceleration
contract is
\[
 a_{min}=\frac{B_{min}-T_{max}-A_{max}}{m_{current}}>0.
\]
For measured distance $d$, speed $v$, total delay $\tau$, geometry error
$e_g$, and terminal tolerance $e_t$, the terminal-entry guard is
\[
 d>0,\quad v>0,\quad B_{min}>T_{max}+A_{max},\quad
 d\geq v\tau+A_{max}\tau^2
       +\frac{v^2}{2a_{min}}+e_g+e_t.
\]
The first term is observation/command/pressure delay distance, the second is
adverse delayed motion, the third is the certified stopping distance, and the
last two are geometry and terminal budgets.  The guard is evaluated at the
current measured speed; there is no separate universal operational speed
theorem.

Lean proves guard positivity, the discrete stopping envelope, and one-step
no-crossing under explicit hypotheses.  The envelope is
\[
 0\leq d,\quad 0\leq v,\quad v^2\leq 2a_{min}d.
\]
It is not enough to prove no crossing.  Arrival is a separate runtime progress
and stable-sample contract.  The mathematical terminal state is $v_s=0$;
runtime acceptance additionally requires
\[
 |v_s|\leq\varepsilon_v,
 \qquad |x|\leq\varepsilon_{stop},
 \qquad \text{stable for }N_{stable}\text{ observations}.
\]
The position tolerance must satisfy
\[
 \varepsilon_{stop}\geq
 \varepsilon_{trace}+\varepsilon_{fit}+\varepsilon_{coupler}
 +\varepsilon_{sample-age}+\varepsilon_{command-delay}
 +\varepsilon_{numeric}.
\]
No decimal tolerance is accepted without a measured or formally bounded term.

\section{Arrival, holding, and failure states}

The proof separates the following claims:
\begin{enumerate}
  \item stopping-envelope preservation;
  \item no-crossing safety;
  \item velocity convergence under the certified discrete brake model;
  \item terminal arrival under a non-vacuous progress witness;
  \item stable terminal regulation and holding;
  \item active-route feasibility and replan behavior.
\end{enumerate}
Lean defines terminal acceptance and progress predicates but does not claim a
terminal-progress theorem from those predicates alone.  This prevents the old
vacuous liveness claim.  The Lua implementation must provide the low-speed IR
progress/holding refinement before any stronger readiness status is possible.

Fail-closed states are required for stale or invalid sample, invalid spline
localization, wrong route orientation, route revision change, missing route,
marker mismatch, unavailable radio/detector data, profile mismatch, invalid
numeric value, pressure delay overrun, and inconsistent plant response.  Each
must emit an in-game message, stop the controller, and require player
resolution.

\section{Claim inventory}

\begin{longtable}{p{0.29\linewidth}p{0.23\linewidth}p{0.39\linewidth}}
\toprule
Claim & Classification & Current evidence \\
\midrule
Bezier endpoint/reverse identities & proved by Lean & \Lean{IRProof} \\
Centre estimate residual implication & proved by Lean under residual certificate & \Lean{TraceCentreSound}; offset and bound are model/runtime inputs \\
IR turn is cubic approximation & derived from bytecode & root \texttt{CubicCurve.circle()} \\
Crossing/overpass routing exclusion & proved by Lean under topology evidence & geometry evidence must be supplied \\
Directed marker mapping & proved by Lean & explicit tolerance and approach \\
Detector order and profile mismatch & proved by Lean under event/info contract & API event delivery is runtime \\
Material constants and default multiplier & derived from bytecode & root classes and configuration \\
Signed force ledger & proved by Lean/derived from bytecode & full Java branch refinement outstanding \\
Combined throttle/brake transition & proved by Lean at contract level & Lua command policy and delay refinement outstanding \\
Guard positivity and one-step envelope & proved by Lean & lower-bound premises remain runtime obligations \\
Uniform IR robust arrival & unsupported/withdrawn & requires full multi-particle, delay, float, and Lua refinement \\
Formal optimality & unsupported/withdrawn & no independent efficiency objective proved \\
\bottomrule
\end{longtable}

\section{Implementation funnel}

The future Lua implementation must add, in order:
\begin{enumerate}
  \item an immutable trace record and compact-tracer centre estimator with
    measured residual bound;
  \item cubic spline storage, arc marks, reverse orientation, incremental
    trace matching, topology ambiguity states, and persistent IDs;
  \item graph rebuild from validated spline intervals and directed marker
    mappings;
  \item detector event catalogue with immediate \texttt{info()}, order,
    UUID, and immutable profile hash;
  \item front-coupler localization from stock definition/bogey/coupler
    offsets, never raw entity position;
  \item route-local validation and revision-triggered replan/fail-closed
    behavior;
  \item command allocation that sends throttle, train brake, independent
    brake, and emergency state together, with delay/pressure-age accounting;
  \item a measured lower-force/upper-traction certificate for every profiled
    stock and curved/slack/push-pull consist;
  \item terminal speed/position deadbands derived from the complete error sum,
    plus stable-observation holding;
  \item adversarial and mutation tests before user approval.
\end{enumerate}

Production Lua is intentionally unchanged by this artifact.  The exact API
calls, state fields, failure messages, and acceptance tests are listed in
\texttt{pid-implementation-plan.md}.

\section{Readiness}

Gate A's Lean portion passes, but the complete Gates A--I do not.  In
particular, the Java-to-Lua floating-point conversion, complete IR particle
linkage/collision refinement, route-local live graph mutation, detector timing,
terminal holding, and robust arrival vector suite are not yet certified by the
formal model.  Therefore the only honest status is
\[
  \boxed{\textbf{NOT\_READY}}.
\]
The final readiness report names each outstanding gate and does not convert a
passing Lean build into implementation approval.

\end{document}
